% !TEX program = xelatex
% AI-effects 프로젝트 --- 한국 직업별 AI 대화량과 SOC AI 지표의 결합, KSCO 8차 변환과 특이사항.
% 표 본문: results/table/23_*.tex (code/04_analysis_kor_ksco8_merge_summary.do가 생성)
% 결합 파일: code/03_merge_kor_conv_ai_ksco8.do → data/proc/merged/kor_{soc6,ksco8}_conv_ai_indicators
% 관련: paper 19(한국 대화량), paper 20(observed exposure 매칭), paper 21(연계표), report 22(연계표 비교)
% 컴파일: xelatex (kotex 필요)
\documentclass[11pt]{article}

\usepackage{fontspec}
\usepackage{kotex}
\usepackage[a4paper,margin=2.2cm]{geometry}
\usepackage{longtable}
\usepackage{booktabs}
\usepackage{array}
\usepackage{amsmath}
\usepackage{xcolor}
\usepackage{hyperref}
\usepackage{enumitem}

\hypersetup{colorlinks=true,linkcolor=blue!50!black,urlcolor=blue!50!black}
\newcolumntype{L}[1]{>{\raggedright\arraybackslash}p{#1}}
\newcolumntype{R}[1]{>{\raggedleft\arraybackslash}p{#1}}
\setlength{\LTleft}{0pt}
\setlength{\LTright}{0pt}
\renewcommand{\arraystretch}{1.15}
\setlength{\tabcolsep}{4pt}

\title{한국 AI 대화량$\cdot$AI 지표의 KSCO 8차 결합\\[2mm]
\large --- 결합 절차와 연계표 결합에서 생긴 특이사항 ---}
\author{AI-effects 프로젝트}
\date{\today}

\begin{document}
\maketitle

\begin{abstract}
\noindent 한국 직업별 월간 업무용 AI 대화량(paper 19)과 SOC 2018 세부 직업별 AI 지표 4종(시간 절감, observed exposure, Eloundou 노출도, Felten AIOE)을 먼저 SOC 단위에서 결합했다. 이어 이 저장소의 SOC 2018–KSCO 8차 연계표(paper 21)로 KSCO 세분류 494개(KECO 2025 포함)로 옮겼다. 대화량은 배분 가중치(\texttt{w\_alloc})로 나누었고, AI 지표는 평균 가중치(\texttt{w\_mean})로 평균했다. 월별 KSCO 대화량의 합은 SOC 대화량의 합과 같다(do 파일에서 확인). 연계표 결합에서 생긴 주요 특이사항은 다섯 가지다.
\begin{itemize}[leftmargin=*]
  \item \textbf{미공개 대화량}: KSCO의 대화량은 연결된 SOC 중 공개된 것만 더한 값이다. 2026년 5월 KSCO 157개는 공개된 SOC가 하나도 없어 결측이다. KSCO 228개(업무용 대화의 44.3\%)는 일부 SOC만 공개되어 대화량이 하한값이다.
  \item \textbf{값이 같은 KSCO}: KSCO 201개(74개 묶음)는 연결된 SOC와 가중치가 같아 대화량과 지표가 모두 같다. 작가와 출판물 전문가, 대학교수$\cdot$대학 강사$\cdot$대학 교육 조교가 그 예다.
  \item \textbf{곱한 양은 SOC에서 계산}: 대화 수 $\times$ 절감 시간 같은 곱은 SOC에서 계산한 뒤 배분해야 한다. 총량은 두 방법이 거의 같지만(100.1\%), KSCO별로는 337개 중 201개에서 10\% 넘게 다르다.
  \item \textbf{``All Other'' 직업의 해석}: 대화량이 큰 SOC 일부는 실제로는 O*NET 세분 직업의 대화다. 예를 들어 15-1299(Web Administrators, 한국 업무용 대화의 9.7\%)가 그렇다. 연계표는 이를 ``All Other'' 범주 전체로 옮긴다.
  \item \textbf{기타}: 4월의 농림어업 자리표시 행(0.04\%)은 배분되지 않는다. 지표 결측은 적다. 대화량이 있는데 observed exposure가 없는 KSCO는 3개(업무용 대화 0.23\%)다.
\end{itemize}
\end{abstract}

\tableofcontents

% ============================================================
\section{목적과 산출 파일}
% ============================================================

이 저장소에는 SOC 2018 세부 직업(6자리) 단위로 만든 자료가 두 가지 있다.
\begin{itemize}[leftmargin=*]
  \item \textbf{한국 업무용 AI 대화량}: \texttt{sensortower\_aei\_kor\_soc6\_conversations.dta}. 월 $\times$ SOC, 2026년 4월 337행, 5월 365행(paper 19).
  \item \textbf{AI 지표}: \texttt{soc6\_ai\_indicators.dta}. SOC 808개에 AEI 시간 절감, observed exposure, Eloundou 노출도, Felten AIOE를 결합한 파일이다.
\end{itemize}
본 노트는 두 자료를 결합하고 SOC 2018–KSCO 8차 연계표(paper 21)로 한국 직업분류에 옮긴 과정을 기록한다. 특히 연계표 결합에서 생긴 특이사항을 정리한다. 결합은 \texttt{code/03\_merge\_kor\_conv\_ai\_ksco8.do}가 하고, 산출 파일은 \texttt{data/proc/merged/}에 있다.
\begin{enumerate}[leftmargin=*]
  \item \texttt{kor\_soc6\_conv\_ai\_indicators.\{csv,dta\}}: 월 $\times$ SOC. AI 지표가 있는 808개 직업 $\times$ 2개월에 4월 자리표시 행 \texttt{45-XXXX} 1행을 더해 1,617행이다. 대화량이 공개되지 않은 직업은 대화량 변수가 결측이고 \texttt{conv\_pub} = 0이다.
  \item \texttt{kor\_ksco8\_conv\_ai\_indicators.\{csv,dta\}}: 월 $\times$ KSCO. 연계표의 KSCO 494개 $\times$ 2개월 = 988행이다. KECO 2025 코드와 명칭을 함께 둔다.
\end{enumerate}

% ============================================================
\section{SOC 단계 결합}
% ============================================================

매칭 키는 월과 \texttt{soc6}다. AI 지표 파일 808개 직업을 월마다 뼈대로 두고 대화량을 붙였다. 대화량이 없는 직업을 버리지 않은 것은, 이 직업들도 KSCO 쪽 연결에는 들어가기 때문이다(4.1절). 표~\ref{tab:socmerge}는 결합 결과다.

\begin{longtable}{L{6.6cm}R{1.8cm}R{2.2cm}R{1.8cm}R{2.2cm}}
\caption{SOC 단계 결합 결과(비중은 한국 업무용 대화 \texttt{share\_work} 기준)}\label{tab:socmerge}\\
\toprule
 & \multicolumn{2}{c}{2026-04} & \multicolumn{2}{c}{2026-05} \\
 & 직업 수 & 비중(\%) & 직업 수 & 비중(\%) \\
\midrule
% 04_analysis_kor_ksco8_merge_summary.do가 생성. 직접 고치지 말 것.
대화량 공개 SOC 행 & 337 & 100.00 & 365 & 100.00 \\
\quad AI 지표 파일과 매칭 & 336 & 99.96 & 365 & 100.00 \\
\quad 매칭 안 됨(45-XXXX 자리표시) & 1 & 0.04 & 0 & 0.00 \\
\quad\quad 시간 절감(AEI) 있음 & 336 & 99.96 & 365 & 100.00 \\
\quad\quad observed exposure 있음 & 317 & 97.82 & 342 & 98.17 \\
\quad\quad Eloundou 노출도 있음 & 336 & 99.96 & 365 & 100.00 \\
\quad\quad Felten AIOE 있음 & 333 & 98.04 & 362 & 98.00 \\
\midrule
AI 지표 808개 중 대화량 미공개 & 472 & -- & 443 & -- \\

\bottomrule
\end{longtable}

\begin{itemize}[leftmargin=*]
  \item \textbf{대화량 직업은 모두 AI 지표 파일에 있다.} 예외는 4월의 \texttt{45-XXXX} 하나다. 이 행은 농림어업 대분류에 공개된 세부 직업이 없어, 대분류 몫을 한 행에 둔 자리표시다(paper 19). 업무용 대화의 0.04\%에 해당한다.
  \item \textbf{지표 보유}: 시간 절감과 Eloundou 노출도는 대화량 직업 전부에 있다. observed exposure는 5월 365개 중 342개에만 있다(paper 20의 23개 미매칭과 같음). 대화 비중으로는 98.2\%를 덮는다.
  \item \textbf{미공개는 0이 아니다.} AI 지표 808개 중 5월 443개는 한국 대화량이 공개되지 않았다. 이 직업의 대화량은 0으로 두지 않고 결측으로 두었다(paper 20의 유의점).
  \item \textbf{명칭 차이}: 두 파일의 직업 명칭이 다른 SOC가 5월에 21개 있다. 그중 20개는 AEI가 .00 코드를 공개하지 않아, O*NET 세분 직업의 이름이 AEI 명칭(\texttt{soc6\_title\_aei})으로 들어온 경우다. 예를 들어 15-1299는 SOC 명칭이 Computer Occupations, All Other이고, AEI 명칭은 Web Administrators다. 이 문제는 4.5절에서 다시 다룬다.
\end{itemize}

% ============================================================
\section{KSCO 변환 규칙}
% ============================================================

SOC 2018 코드를 $s$, KSCO를 $k$라 하자. 가중치는 paper 21의 연계표에서 가져온다.
\begin{itemize}[leftmargin=*]
  \item \textbf{양(대화 수 \texttt{conv\_k*}, 비중 \texttt{share\_work}, \texttt{pct})}: $Y_k = \sum_s w^{\text{alloc}}_{sk} Y_s$. 대화량이 공개된 SOC만 더한다. 연결된 SOC 가운데 공개된 것이 없으면 결측이다.
  \item \textbf{비율 지표(시간 절감, observed exposure, Eloundou, Felten)}: $x_k = \sum_s w^{\text{mean}}_{sk} x_s / \sum_s w^{\text{mean}}_{sk}$. 값이 있는 SOC만 쓰고, 그 SOC의 가중치 합으로 다시 나눈다.
  \item \textbf{덮개 변수}: 값이 있는 SOC의 $w^{\text{mean}}$ 합을 \texttt{cov\_aei}, \texttt{cov\_obs}, \texttt{cov\_elo}, \texttt{cov\_felten}으로 둔다. 대화량이 공개된 SOC의 $w^{\text{mean}}$ 합은 \texttt{conv\_cover}다. 모두 0~1이고, 1이면 연결된 SOC 전부에 값이 있다.
  \item \textbf{보조 변수}: 연결된 SOC 수는 \texttt{n\_soc\_link}, 그중 대화량이 공개된 SOC 수는 \texttt{n\_soc\_conv}다.
\end{itemize}
AI 지표는 월과 무관하다. 시간 절감은 2026년 4$\cdot$5월을 합친 값이다. 그래서 두 달에 같은 값이 들어간다. 또한 시간 절감은 AEI의 전 세계(GLOBAL) 값이다. 세부 직업 단위의 시간 지표는 국가별로 공개되지 않기 때문이다. 이 값은 한국 대화량과 짝지어 쓰지만 한국 값은 아니다.

% ============================================================
\section{연계표 결합에서 생긴 특이사항}
% ============================================================

\subsection{연계표에는 있으나 AI 지표가 없는 SOC 59개}

연계표의 SOC 2018은 867개이고, AI 지표 파일은 808개다. 나머지 59개는 어느 지표 자료에도 없는 SOC다(paper 21 표 15). 이 59개 SOC가 연결된 쌍(169쌍)도 KSCO 쪽에 남겼다. 그래서 \texttt{n\_soc\_link}가 연계표의 SOC 수(\texttt{n18\_perk})와 같다. 이 SOC들은 값이 없으므로 지표 평균에서는 빠지고, 덮개 변수를 1보다 작게 만든다.

\subsection{4월 자리표시 행은 배분되지 않는다}

\texttt{45-XXXX}는 SOC 세부 직업 코드가 아니어서 연계표에 없다. 따라서 2026년 4월 KSCO 대화량의 합은 업무용 대화의 99.96\%다(표~\ref{tab:kscoconv}의 합계). 이 몫을 농림어업 관련 KSCO에 고르게 나누는 방법도 있으나, 대상 직업을 정할 근거가 없어 하지 않았다.

\subsection{대화량 공개 범위: 결측과 하한값}

\begin{longtable}{L{4.6cm}R{1.2cm}R{1.8cm}R{1.8cm}R{1.2cm}R{1.8cm}R{1.8cm}}
\caption{KSCO 494개의 대화량 공개 범위}\label{tab:kscoconv}\\
\toprule
 & \multicolumn{3}{c}{2026-04} & \multicolumn{3}{c}{2026-05} \\
연결된 SOC의 대화량 공개 & KSCO 수 & 대화 비중(\%) & 평균 \texttt{conv\_cover} & KSCO 수 & 대화 비중(\%) & 평균 \texttt{conv\_cover} \\
\midrule
% 04_analysis_kor_ksco8_merge_summary.do가 생성. 직접 고치지 말 것.
연결된 SOC 모두 공개 & 99 & 53.16 & 1.000 & 109 & 55.71 & 1.000 \\
일부만 공개 & 213 & 46.80 & 0.453 & 228 & 44.29 & 0.470 \\
공개된 SOC 없음(대화량 결측) & 182 & -- & 0.000 & 157 & -- & 0.000 \\
\midrule
합계 & 494 & 99.96 & 0.396 & 494 & 100.00 & 0.438 \\

\bottomrule
\end{longtable}

\begin{itemize}[leftmargin=*]
  \item \textbf{대화량 결측 KSCO}: 5월 기준 157개(4월 182개)다. 연결된 SOC가 모두 미공개이므로, 대화량이 ``공개 기준 미만''이라는 뜻이지 0이라는 뜻이 아니다. 157개 중 93개가 KSCO 대분류 7(기능원)과 8(장치$\cdot$기계 조작 및 조립 종사자)이고, 군인 5개도 모두 여기에 든다.
  \item \textbf{일부만 공개된 KSCO}: 5월 기준 228개다. 업무용 대화의 44.3\%가 이 KSCO들에 들어간다. 대화량은 공개된 SOC 몫만 더한 값이므로 하한값이다. 다만 미공개 SOC의 몫은 국가 전체로도 1.92\% 안이라(paper 20), 과소 정도는 작을 것으로 보인다.
  \item 표~\ref{tab:partial}은 일부 공개 KSCO 가운데 대화량이 큰 10개다. 통계$\cdot$데이터 관련 사무원(3320)은 연결된 SOC 12개 중 8개, 자연과학 시험원(2133)은 18개 중 11개만 공개되었다.
\end{itemize}

\begin{longtable}{L{8.4cm}R{2.4cm}R{2.6cm}R{2.4cm}}
\caption{2026-05 대화량 일부 공개 KSCO 중 대화량 상위 10개}\label{tab:partial}\\
\toprule
KSCO & 대화 비중(\%) & 공개 SOC / 연결 SOC & \texttt{conv\_cover} \\
\midrule
% 04_analysis_kor_ksco8_merge_summary.do가 생성. 직접 고치지 말 것.
3320 통계‧데이터 관련 사무원 & 2.67 & 8/12 & 0.833 \\
2133 자연과학 시험원 & 2.51 & 11/18 & 0.800 \\
2914 번역가 및 통역가 & 1.96 & 1/2 & 0.500 \\
2239 기타 네트워크 및 정보 보안 전문가 & 1.86 & 5/6 & 0.500 \\
2232 정보 보안 전문가 & 1.86 & 5/6 & 0.500 \\
2933 아나운서 및 리포터 & 1.65 & 2/3 & 0.667 \\
3710 총무 사무원 및 대학 행정조교 & 1.33 & 2/3 & 0.750 \\
3721 비서 & 1.31 & 4/6 & 0.750 \\
3409 기타 금융 사무원 & 1.25 & 4/8 & 0.667 \\
3403 은행 사무원 & 1.25 & 4/8 & 0.667 \\

\bottomrule
\end{longtable}

\subsection{값이 같아지는 KSCO}

연계표에서 연결된 SOC 집합과 가중치(\texttt{w\_mean}, \texttt{w\_alloc})가 같은 KSCO는 대화량과 모든 지표가 같은 값을 받는다. 이런 KSCO가 201개(74개 묶음)로, KSCO 494개의 41\%다. 표~\ref{tab:twins}는 대화량이 큰 묶음이다.

\begin{longtable}{L{11.4cm}R{1.6cm}R{2.8cm}}
\caption{값이 같아지는 KSCO 묶음(대화량 상위 12개, 비중은 KSCO 하나당 2026-05 업무용 대화 비중)}\label{tab:twins}\\
\toprule
KSCO 묶음 & KSCO 수 & 대화 비중(\%) \\
\midrule
\endfirsthead
\toprule
KSCO 묶음 & KSCO 수 & 대화 비중(\%) \\
\midrule
\endhead
% 04_analysis_kor_ksco8_merge_summary.do가 생성. 직접 고치지 말 것.
2911 작가; 2912 출판물 전문가 & 2 & 3.38 \\
2232 정보 보안 전문가; 2239 기타 네트워크 및 정보 보안 전문가 & 2 & 1.86 \\
3403 은행 사무원; 3409 기타 금융 사무원 & 2 & 1.25 \\
2831 상품 기획 전문가; 2832 여행 상품 개발자 & 2 & 0.95 \\
2611 대학교수; 2612 대학 강사; 2692 대학 교육 조교 & 3 & 0.48 \\
2241 데이터 시스템 전문가; 2242 데이터 분석가 & 2 & 0.54 \\
2211 컴퓨터 하드웨어 기술자 및 연구원; 2342 반도체공학 기술자 및 연구원; 2343 전자공학 기술자 및 연구원 & 3 & 0.36 \\
2981 공연 및 시각예술 기획자; 2982 영화 및 음반 기획자; 3722 연예인 및 스포츠 매니저 & 3 & 0.32 \\
2849 기타 기술 영업 및 중개 관련 종사원; 5101 자동차 영업원 & 2 & 0.25 \\
2375 신재생에너지 시험원; 2376 가스 및 에너지 시험원; 2393 식품공학 시험원; 2394 섬유공학 시험원 & 4 & 0.11 \\
2352 로봇공학 기술자 및 연구원; 2361 방재 기술자 및 연구원; 2392 섬유공학 기술자 및 연구원 & 3 & 0.13 \\
9923 주차 관리원 및 안내원; 9993 검표원; 9994 대여 제품 방문 점검원 & 3 & 0.12 \\
\midrule
합계(묶음 74개) & 201 & -- \\

\bottomrule
\end{longtable}

\begin{longtable}{L{11.0cm}R{2.0cm}R{2.4cm}}
\caption{값이 같아지는 KSCO 묶음이 생기는 구조}\label{tab:twinstruct}\\
\toprule
 & 수 & 비율(\%) \\
\midrule
% 04_analysis_kor_ksco8_merge_summary.do가 생성. 직접 고치지 말 것.
값이 같아지는 KSCO & 201 & 100.0 \\
\quad ISCO 1개에서만 값을 받는 KSCO & 188 & 93.5 \\
\quad SOC 1개만 받는 KSCO & 28 & 13.9 \\
\quad SOC 2개 이상 받는 KSCO(최대 37개) & 173 & 86.1 \\
\midrule
묶음 & 74 & 100.0 \\
\quad 구성원이 받는 ISCO 집합이 모두 같은 묶음 & 72 & 97.3 \\

\bottomrule
\end{longtable}

\begin{itemize}[leftmargin=*]
  \item \textbf{원인}: ISCO 하나가 KSCO 여럿으로 나뉘고(3단계 1:N), 나뉜 KSCO들이 다른 ISCO를 받지 않는 경우다. 그러면 그 ISCO로 모인 SOC들의 값이 KSCO마다 똑같이 복사된다(표~\ref{tab:twinstruct}). 묶음 KSCO 201개 가운데 188개(93.5\%)는 ISCO 1개에서만 값을 받는다. 74개 묶음 가운데 72개(97.3\%)는 구성원이 받는 ISCO 집합이 모두 같다. 예를 들어 ISCO 2641(Authors and related writers)은 작가(2911)와 출판물 전문가(2912)로 나뉜다. 두 KSCO는 같은 SOC들의 값을 같은 비율로 받는다. 대학교수$\cdot$대학 강사$\cdot$대학 교육 조교는 ISCO 2310 하나를 나눠 받는다. ISCO 2310에는 대학 교원 SOC 37개가 모이므로, 세 KSCO가 같은 SOC 37개의 값을 받는다.
  \item \textbf{``SOC 하나가 KSCO 여럿과 연결''과는 다르다}: 묶음 안의 SOC는 모두 묶음 안의 KSCO 전부와 연결되므로, SOC 쪽에서 보면 1:N이다. 그러나 KSCO 2개 이상과 연결되는 SOC는 516개(paper 21 표 4)로 훨씬 많고, 대부분은 묶음을 만들지 않는다. 받는 KSCO가 다른 ISCO나 SOC도 함께 받아 SOC 집합이 달라지기 때문이다. 예를 들어 Software Developers(15-1252)는 시스템 소프트웨어 개발자(2222)와 응용 소프트웨어 개발자(2223)에 함께 연결된다. 그러나 2223은 ISCO 2513$\cdot$2514를 거쳐 다른 SOC도 받으므로 두 KSCO의 값이 다르다. 또한 묶음 KSCO의 86.1\%는 SOC를 2개 이상(최대 37개) 받는다. SOC 1개만 받는 KSCO는 28개(13.9\%)뿐이다.
  \item \textbf{의미가 먼 묶음}: 상품 기획 전문가와 여행 상품 개발자는 ISCO 2431(광고$\cdot$마케팅 전문가)을 나눠 받는다. 그래서 Writers and Authors의 대화가 여행 상품 개발자에게도 같은 크기로 간다(paper 21, 8.1절).
  \item \textbf{분석상 유의점}: 같은 묶음 안의 KSCO는 AI 지표에 변이가 없다. KSCO 단위 회귀에서는 묶음 내부의 고용$\cdot$임금 차이가 AI 지표로 설명될 수 없다. 표준오차는 묶음(또는 ISCO) 단위로 군집화하는 것을 검토해야 한다. KSCO 단위 결과를 보고할 때 묶음 안의 순위를 해석해서는 안 된다.
\end{itemize}

\subsection{``All Other'' 직업과 O*NET 세분 직업}

AEI의 직업 단위는 O*NET-SOC 8자리다. .00 코드가 공개되지 않고 세분 직업(.01 이상)만 공개된 SOC는, 그 세분 직업의 대화가 SOC 전체의 대화로 합산된다(paper 20). 한국 대화량에서 이런 SOC가 20개다. 이 가운데 대화량이 큰 것은 대부분 ``All Other'' 범주다. 대표적인 예는 15-1299(Computer Occupations, All Other, 실제 대화는 Web Administrators, 업무용 대화의 9.7\%)다. 이 밖에 17-2199(Energy Engineers 등), 19-1029(Bioinformatics Scientists), 13-1199(Business Continuity Planners), 19-4099(Quality Control Analysts), 15-2099(Bioinformatics Technicians)가 있다. 연계표는 이 대화를 ``All Other'' 범주 전체의 ISCO로 보낸다. 그래서 실제 과업과 KSCO 배분이 어긋날 수 있다. 15-1299의 대화 절반이 정보 시스템 운영자(2251)로 가서 이 KSCO가 한국 업무용 대화 1위(7.25\%)가 되는 것이 그 예다(표~\ref{tab:top}).

\begin{longtable}{L{6.4cm}R{1.6cm}R{1.8cm}R{1.8cm}R{1.8cm}R{2.2cm}}
\caption{2026-05 업무용 대화 상위 15개 KSCO}\label{tab:top}\\
\toprule
KSCO & 대화 비중(\%) & 대화 수(만 건, $k$ = 3) & 공개 / 연결 SOC & 시간 절감(분) & observed exposure \\
\midrule
% 04_analysis_kor_ksco8_merge_summary.do가 생성. 직접 고치지 말 것.
2251 정보 시스템 운영자 & 7.25 & 2,512 & 3/3 & 124 & 0.355 \\
2229 기타 컴퓨터 시스템 및 소프트웨어 전문가 & 3.73 & 1,291 & 5/5 & 336 & 0.414 \\
2911 작가 & 3.38 & 1,169 & 3/3 & 251 & 0.322 \\
2912 출판물 전문가 & 3.38 & 1,169 & 3/3 & 251 & 0.322 \\
2913 기자 및 언론 관련 전문가 & 3.34 & 1,156 & 2/2 & 247 & 0.228 \\
2922 사서 및 기록물 관리사 & 3.17 & 1,097 & 1/1 & 122 & 0.203 \\
2223 응용 소프트웨어 개발자 & 2.85 & 988 & 5/5 & 423 & 0.504 \\
3320 통계‧데이터 관련 사무원 & 2.67 & 925 & 8/12 & 377 & 0.299 \\
2133 자연과학 시험원 & 2.51 & 868 & 11/18 & 307 & 0.253 \\
2221 컴퓨터 시스템 전문가 & 2.21 & 766 & 3/3 & 241 & 0.323 \\
2914 번역가 및 통역가 & 1.96 & 678 & 1/2 & 284 & 0.232 \\
2239 기타 네트워크 및 정보 보안 전문가 & 1.86 & 646 & 5/6 & 287 & 0.454 \\
2232 정보 보안 전문가 & 1.86 & 646 & 5/6 & 287 & 0.454 \\
2933 아나운서 및 리포터 & 1.65 & 571 & 2/3 & 317 & 0.137 \\
3724 속기사 & 1.57 & 543 & 1/1 & 258 & 0.238 \\

\bottomrule
\end{longtable}

\noindent 상위 KSCO에는 자연과학 시험원(2133, 2.51\%)도 있다. 이 대화의 69\%는 ISCO 3314(통계$\cdot$수학 준전문가)를 거쳐 들어온다. Data Scientists(15-2051)와 Statistical Assistants(43-9111) 등의 대화가 그 경로를 탄다. 통계청 공식 연계(3314 → 2133, 3320)를 따른 결과지만, 의미상 통계$\cdot$데이터 관련 사무원(3320)에 더 가깝다. 이 연계는 report 22에서 외부 연계표가 공백 때문에 빠뜨린 바로 그 연계다. 어느 쪽을 쓰느냐에 따라 2133의 대화량이 크게 달라진다.

\subsection{AI 지표의 덮개와 결측}

\begin{longtable}{L{3.6cm}R{1.4cm}R{1.4cm}R{1.6cm}R{1.6cm}R{1.6cm}R{1.4cm}R{2.4cm}}
\caption{지표별 SOC$\cdot$KSCO 보유와 덮개(KSCO는 2026-05)}\label{tab:indcov}\\
\toprule
지표 & SOC 수 & KSCO 수 & 덮개 = 1 & 덮개 $<$ 1 & 평균 덮개 & 결측 KSCO & 결측 KSCO의 대화 비중(\%) \\
\midrule
% 04_analysis_kor_ksco8_merge_summary.do가 생성. 직접 고치지 말 것.
시간 절감(AEI, 분) & 613 & 455 & 230 & 225 & 0.801 & 39 & 0.00 \\
observed exposure & 756 & 480 & 302 & 178 & 0.901 & 14 & 0.23 \\
Eloundou GPT-4 beta & 798 & 488 & 368 & 120 & 0.944 & 6 & 0.00 \\
Felten 언어모델 AIOE & 800 & 484 & 384 & 100 & 0.951 & 10 & 0.00 \\

\bottomrule
\end{longtable}

\begin{itemize}[leftmargin=*]
  \item \textbf{시간 절감}: SOC 613개에만 있어 KSCO 39개가 결측이다. 그러나 결측 KSCO의 한국 대화 비중은 0이다. 한국 대화량이 공개된 SOC는 모두 전 세계 시간 지표도 공개되어 있기 때문이다(표~
ef{tab:socmerge}). 대신 KSCO 225개는 덮개가 1보다 작다(평균 0.80). 이 KSCO의 시간 절감은 연결된 SOC 일부의 평균이다.
  \item \textbf{observed exposure}: KSCO 14개가 결측이다. 이 가운데 대화량이 있는 KSCO는 3개로, 업무용 대화의 0.23\%다(표~\ref{tab:missind}).
  \item \textbf{Eloundou, Felten}: 덮개가 높고(평균 0.94~0.95) 결측 KSCO의 대화 비중도 0이다.
\end{itemize}

\begin{longtable}{L{8.4cm}R{2.4cm}R{2.6cm}R{2.4cm}}
\caption{2026-05 대화량은 있으나 observed exposure가 없는 KSCO}\label{tab:missind}\\
\toprule
KSCO & 대화 비중(\%) & 공개 / 연결 SOC & 시간 절감(분) \\
\midrule
% 04_analysis_kor_ksco8_merge_summary.do가 생성. 직접 고치지 말 것.
2842 해외 영업원 & 0.180 & 2/3 & 269 \\
4211 교사 보조 및 관련 종사원 & 0.050 & 2/3 & 255 \\
9991 구두 미화원 & 0.000 & 1/2 & 76 \\

\bottomrule
\end{longtable}

\subsection{곱한 양은 SOC 단계에서 계산해야 한다}

``대화 수 $\times$ 대화당 절감 시간''처럼 양과 비율을 곱한 총량은 SOC에서 곱한 뒤 \texttt{w\_alloc}으로 배분해야 한다(A). KSCO 파일의 배분된 대화 수와 KSCO 평균 절감 시간을 곱하면(B) 다른 값이 나온다. 표~\ref{tab:product}는 2026년 5월 업무용 대화($k$ = 3)로 두 방법을 비교한 것이다.

\begin{longtable}{L{10.6cm}R{2.8cm}R{2.4cm}}
\caption{2026-05 총 절감 시간: SOC에서 곱함(A) 대 KSCO에서 곱함(B)}\label{tab:product}\\
\toprule
 & 백만 시간 & (A) 대비(\%) \\
\midrule
% 04_analysis_kor_ksco8_merge_summary.do가 생성. 직접 고치지 말 것.
(A) SOC에서 곱한 뒤 \texttt{w\_alloc}으로 배분해 합산 & 1499.89 & 100.0 \\
(B) KSCO 대화 수 × KSCO 평균 절감 시간(\texttt{w\_mean}) & 1500.78 & 100.1 \\
\midrule
KSCO별 (A)와 (B)의 상관계수 & 0.983 & -- \\
KSCO별 차이가 10\% 넘는 KSCO 수 / 값이 있는 KSCO 수 & 201 / 337 & -- \\

\bottomrule
\end{longtable}

\begin{itemize}[leftmargin=*]
  \item \textbf{총량}: 두 방법의 차이는 0.1\%로 작다. 그러나 KSCO별로는 337개 중 201개에서 10\% 넘게 다르다. 차이의 중앙값은 0이고, 사분위 범위는 $-14$\%~$+15$\%, 최대는 $+269$\%다.
  \item \textbf{원인}: (B)는 KSCO 평균 절감 시간을 대화량과 무관하게 같은 가중치(\texttt{w\_mean})로 만든다. 대화가 많은 SOC와 적은 SOC의 절감 시간이 같은 무게로 섞이고, 대화량은 \texttt{w\_alloc}으로 따로 배분된다. 두 가중치가 다르므로 곱의 합이 맞지 않는다.
  \item \textbf{권고}: 총 절감 시간, 노출도 가중 대화량 같은 곱은 SOC 파일(\texttt{kor\_soc6\_conv\_ai\_indicators})에서 계산한 뒤 연계표의 \texttt{w\_alloc}으로 KSCO에 배분한다. KSCO 파일의 지표 평균은 ``그 KSCO에 연결된 미국 직업들의 평균적 특성''으로만 쓴다.
\end{itemize}

\subsection{KECO 2025}

KSCO와 KECO는 1:1이므로, KECO 기준 분석에는 \texttt{keco2025}를 키로 같은 파일을 쓰면 된다. KECO 8352(KSCO 8631, 일차전지 및 이차전지 제조 기계 조작원)는 연계표에 없어 이 파일에도 없다(paper 21, 6.3절).

% ============================================================
\section{정리}
% ============================================================

\begin{enumerate}[leftmargin=*]
  \item SOC 단계 결합은 문제가 거의 없다. 대화량 직업은 모두 AI 지표 파일에 있고, 지표 결측도 대화 비중으로 2\% 미만이다.
  \item 연계표 결합에서는 1:1이 아닌 대응 때문에 네 가지를 주의해야 한다. 첫째, 대화량이 결측이거나 하한값인 KSCO가 있다. 둘째, KSCO의 41\%가 다른 KSCO와 같은 값을 받는다. 셋째, 곱한 양은 SOC에서 계산해야 한다. 넷째, ``All Other''와 O*NET 세분 직업의 대화가 넓게 배분된다.
  \item 대화량 상위 KSCO는 연계 경로에 민감하다. 정보 시스템 운영자(15-1299 경로)와 자연과학 시험원(ISCO 3314 경로)이 대표적이다. KSCO 단위 결과를 해석할 때는 \texttt{soc2018\_ksco8\_crosswalk}의 \texttt{via\_soc2010}, \texttt{via\_isco08}로 경로를 확인한다.
\end{enumerate}

% ============================================================
\section{재현}
% ============================================================

Stata 17에서 다음 순서로 실행한다. 1번의 세 파일은 기존 결과를 만드는 do 파일이다.
\begin{enumerate}[leftmargin=*]
  \item \texttt{03\_merge\_sensortower\_aei\_kor\_soc\_conversations.do}, \texttt{03\_merge\_soc6\_ai\_indicators.do}, \texttt{01\_import\_soc2018\_ksco8\_crosswalk.do}
  \item \texttt{03\_merge\_kor\_conv\_ai\_ksco8.do}: 결합 파일 2종을 만든다.
  \item \texttt{04\_analysis\_kor\_ksco8\_merge\_summary.do}: 본 노트의 표(\texttt{results/table/23\_*.tex})와 \texttt{23\_kor\_ksco8\_conv\_ai.xlsx}(5월 KSCO별 전체 변수)를 만든다.
\end{enumerate}

\end{document}
